\chapter{Der Stand im Jahr 2026}\label{kap:aktuell}

\begin{merksatz}
\boxtitel{Was gilt und was behauptet wird}
\begin{center}
{\small
\renewcommand{\arraystretch}{1.5}
\begin{tabular}{@{}P{3.6cm}P{2.0cm}P{2.4cm}@{}}
\toprule
\textbf{Aussage} & \textbf{Beleglage} & \textbf{Urteil}\\
\midrule
Fehler sitzen an den Nähten & \textcolor{gruen}{gut} & \textcolor{gruen}{$\approx45\,\%$}\\
Der Mittelwert taugt nicht & \textcolor{gruen}{eindeutig} & \textcolor{gruen}{Kap.~\ref{kap:messen}}\\
Idempotenz ist unverzichtbar & \textcolor{gruen}{eindeutig} & \\
Verfügbarkeit multipliziert sich & \textcolor{gruen}{\textbf{Rechnung}} & \textcolor{gruen}{\textbf{Kap.~\ref{kap:teilen}}}\\
Beobachtbarkeit gehört zum Entwurf & \textcolor{gruen}{gut} & \\
\midrule
\enquote{Aufteilen macht sauberer} & \textcolor{akzent}{\textbf{falsch}} & \textcolor{akzent}{\textbf{Module genügen}}\\
\enquote{Verteilt ist skalierbar} & \textcolor{akzent}{unvoll\-ständig} & \textcolor{akzent}{erst messen}\\
\enquote{Ein Format für alles} & \textcolor{akzent}{falsch} & \textbf{eines je Naht}\\
\bottomrule
\end{tabular}}
\end{center}

\begin{center}
\begin{minipage}{0.88\textwidth}
\kursiv{Die vierte Zeile ist keine Erfahrung, sondern eine Rechnung --- und deshalb die belastbarste Aussage des Bandes. Wer die Zahl der Bestandteile erhöht, senkt die Verfügbarkeit, es sei denn, er erhöht die jedes einzelnen.}
\end{minipage}
\end{center}
\end{merksatz}

\begin{aktuell}
\boxtitel{Fünf Entwicklungen}
\textbf{Erstens: Die Pendelbewegung.}

\begin{center}
\renewcommand{\arraystretch}{1.45}
\begin{tabular}{@{}P{2.2cm}P{2.8cm}P{2.8cm}@{}}
\toprule
\textbf{Zeitraum} & \textbf{Vorherrschend} & \\
\midrule
$2005$ & Große Einzelanwendung & \\
$2015$ & \textbf{möglichst viele Dienste} & \textcolor{akzent}{Übertreibung}\\
$2022$ & Rückbau & \textcolor{gruen}{gemessene Kosten}\\
$2026$ & \textbf{wenige, begründete} & \textcolor{gruen}{\textbf{nach Anforderung}}\\
\bottomrule
\end{tabular}
\end{center}

\begin{center}
\begin{minipage}{0.86\textwidth}
\kursiv{Bemerkenswert ist die dritte Zeile: Mehrere große Anbieter haben Dienste wieder zusammengelegt und die Kosten offengelegt. Das hat die Diskussion versachlicht --- die Frage lautet nicht mehr \enquote{wie viele Dienste}, sondern \enquote{welcher Grund}.}
\end{minipage}
\end{center}

\medskip
\textbf{Zweitens: Fachstandards setzen sich durch.}

\begin{center}
\renewcommand{\arraystretch}{1.45}
\begin{tabular}{@{}P{3.0cm}P{2.2cm}P{2.6cm}@{}}
\toprule
\textbf{Bereich} & \textbf{Stand} & \\
\midrule
Befundaustausch & \textcolor{gruen}{\textbf{etabliert}} & \textbf{FHIR}\\
Codelisten & \textcolor{gruen}{etabliert} & LOINC, UCUM\\
Verordnung & \textcolor{gruen}{in Umsetzung} & \textcolor{gruen}{EU}\\
Hausinterne Nähte & \textcolor{mittel}{eigen} & \textcolor{gruen}{zulässig}\\
\bottomrule
\end{tabular}
\end{center}

\begin{center}
\kursiv{Die letzte Zeile ist die praktische Empfehlung aus Kapitel~\ref{kap:vertraege}: Fachstandard nach außen, eigenes Format innen, Übersetzung an einer Stelle. Das ist inzwischen die verbreitete Bauform und regulatorisch anerkannt.}
\end{center}

\medskip
\textbf{Drittens: Beobachtbarkeit ist Standard geworden.}

\begin{center}
\renewcommand{\arraystretch}{1.45}
\begin{tabular}{@{}P{3.0cm}P{2.2cm}P{2.6cm}@{}}
\toprule
\textbf{Praxis} & \textbf{Verbreitung} & \\
\midrule
Zählwerte & \textcolor{gruen}{\textbf{üblich}} & \\
Verfolgung über Nähte & \textcolor{gruen}{verbreitet} & \textbf{einheitliches Format}\\
Strukturierte Protokolle & \textcolor{gruen}{üblich} & \\
Ziele schriftlich & \textcolor{mittel}{seltener} & \textcolor{akzent}{\textbf{die Lücke}}\\
\bottomrule
\end{tabular}
\end{center}

\begin{center}
\begin{minipage}{0.86\textwidth}
\kursiv{Die letzte Zeile ist der verbleibende Mangel: Die Werkzeuge sind da, die Zahlen werden erhoben --- und niemand hat festgelegt, welcher Wert gut genug ist. Ohne Zielvorgabe ist ein Schaubild nur ein Schaubild.}
\end{minipage}
\end{center}

\medskip
\textbf{Viertens: Regulierung erreicht die Nähte.}

\begin{center}
\renewcommand{\arraystretch}{1.45}
\begin{tabular}{@{}P{3.0cm}P{2.2cm}P{2.6cm}@{}}
\toprule
\textbf{Anforderung} & \textbf{Betrifft} & \\
\midrule
Nachvollzieh\-barkeit & \textbf{beide Zeitachsen} & \textcolor{gruen}{Kap.~\ref{kap:zeit}}\\
Stückliste & alle Bestandteile & \\
Vollständigkeit belegen & \textbf{Zählwerte} & \textcolor{gruen}{\textbf{Kap.~\ref{kap:beobachtbarkeit}}}\\
Löschfristen & \textbf{auch Zwischenspeicher} & \textcolor{akzent}{Kap.~\ref{kap:zwischenspeicher}}\\
\bottomrule
\end{tabular}
\end{center}

\begin{center}
\kursiv{Die dritte Zeile ist neu und folgenreich: Es genügt nicht, dass Befunde ankommen --- man muss belegen können, dass keiner fehlt. Das verlangt einen Abgleich der Zählwerte über die ganze Kette, und den hat kaum jemand.}
\end{center}

\medskip
\textbf{Fünftens: Werkzeuge, die Nähte erzeugen.}

\begin{center}
\renewcommand{\arraystretch}{1.45}
\begin{tabular}{@{}P{3.0cm}P{2.2cm}P{2.6cm}@{}}
\toprule
\textbf{Wirkung} & \textbf{Befund} & \\
\midrule
Gerüst und Umsetzung & \textcolor{gruen}{\textbf{gut}} & \textbf{aus dem Schema}\\
Fehler\-behandlung & \textcolor{mittel}{brauchbar} & \\
\textbf{Vertragsfragen} & \textcolor{akzent}{\textbf{unbeant\-wortet}} & \textcolor{akzent}{\textbf{Einheit, Zeitpunkt}}\\
Idempotenz & \textcolor{akzent}{selten} & \textcolor{akzent}{muss man verlangen}\\
\bottomrule
\end{tabular}
\end{center}

\begin{center}
\begin{minipage}{0.86\textwidth}
\kursiv{Die dritte Zeile ist der Kern: Ein Werkzeug erzeugt aus einem Schema tadellosen Code --- und das Schema enthält genau die sieben Punkte aus Kapitel~\ref{kap:vertraege} nicht, die den Fall aus dem Vorwort verursacht haben. Es kann nur umsetzen, was jemand entschieden hat.}
\end{minipage}
\end{center}
\end{aktuell}

\begin{merksatz}
\boxtitel{Drei Einsichten}
\textbf{Erstens: Die Naht ist die teuerste Entscheidung.}

\begin{center}
\renewcommand{\arraystretch}{1.45}
\begin{tabular}{@{}P{2.8cm}P{2.2cm}P{2.6cm}@{}}
\toprule
\textbf{Gegenstand} & \textbf{Lebensdauer} & \textbf{Aufwand üblich}\\
\midrule
Code & $\approx2$ Jahre & \textcolor{gruen}{hoch}\\
Bestandteil & $\approx7$ Jahre & \textcolor{gruen}{hoch}\\
\textbf{Naht} & $\approx\mathbf{20}$ \textbf{Jahre} & \textcolor{akzent}{\textbf{gering}}\\
\textbf{Kennungs\-schema} & \textcolor{akzent}{\textbf{unbegrenzt}} & \textcolor{akzent}{\textbf{sehr gering}}\\
\bottomrule
\end{tabular}
\end{center}

\begin{center}
\kursiv{Das Verhältnis ist umgekehrt proportional: Je länger etwas lebt, desto weniger Sorgfalt bekommt es. Das ist der Befund, aus dem dieser Band entstanden ist.}
\end{center}

\textbf{Zweitens: Fast die Hälfte der Fehler ist nicht Code.}

\begin{center}
{\small
\renewcommand{\arraystretch}{1.45}
\begin{tabular}{@{}P{3.0cm}P{2.4cm}P{2.4cm}@{}}
\toprule
\textbf{Fehlerklasse} & \textbf{Ursache} & \textbf{Kapitel}\\
\midrule
Falscher Zeitpunkt & \textbf{Vereinbarung} & \ref{kap:zeit}\\
Falsche Einheit & Vereinbarung & \ref{kap:zahlen}\\
Werte doppelt & \textbf{fehlende Kennung} & \ref{kap:naht}\\
Lücke unsichtbar & \textbf{Anzeige\-entscheidung} & \ref{kap:fehler}\\
Kaskade & fehlende Schranke & \ref{kap:fehler}\\
Veralteter Wert & Zwischen\-speicher & \ref{kap:zwischenspeicher}\\
\bottomrule
\end{tabular}}
\end{center}

\begin{center}
\begin{minipage}{0.88\textwidth}
\kursiv{Sechs Fehlerklassen, keine davon ein Programmierfehler. Jede entsteht dort, wo zwei richtige Programme sich über etwas nicht geeinigt haben --- und keine wird von einem Werkzeug gefunden, das ein Programm prüft.}
\end{minipage}
\end{center}

\textbf{Drittens: Das Handwerk trägt auch hier.}

\begin{center}
\renewcommand{\arraystretch}{1.5}
\begin{tabular}{@{}P{3.0cm}P{2.2cm}P{2.6cm}@{}}
\toprule
\textbf{Aus Band 20} & \textbf{Gilt hier?} & \\
\midrule
Aufgabe verstehen & \textcolor{gruen}{\textbf{unverändert}} & \\
Daten zuerst & \textcolor{gruen}{\textbf{verstärkt}} & \textbf{der Vertrag}\\
Fehlersuche als Verfahren & \textcolor{gruen}{unverändert} & \textbf{plus Verfolgung}\\
Fremden Code lesen & \textcolor{gruen}{unverändert} & \\
Wann aufhören & \textcolor{gruen}{plus drei Kriterien} & \\
\bottomrule
\end{tabular}
\end{center}

\begin{center}
\kursiv{Viermal dasselbe Ergebnis über vier Bände. Die zweite Zeile fällt hier am stärksten aus: Der Vertrag einer Naht ist ein Datenmodell, das zwanzig Jahre gilt --- Wirths Rangordnung auf ihrer längsten Zeitskala.}
\end{center}

\begin{warnung}
\boxtitelr{Was dieser Band nicht behauptet}
\begin{center}
\renewcommand{\arraystretch}{1.5}
\begin{tabular}{@{}P{3.4cm}P{2.2cm}P{2.2cm}@{}}
\toprule
\textbf{Behauptung} & \textbf{Aufgestellt?} & \\
\midrule
Verteilen ist schlecht & \textcolor{akzent}{\textbf{nein}} & \textbf{vier Gründe}\\
Es gibt eine richtige Zahl von Diensten & \textcolor{akzent}{\textbf{nein}} & \\
Nähte lassen sich fehlerfrei bauen & \textcolor{akzent}{nein} & \\
Diese Zahlen sind Konstanten & \textcolor{akzent}{nein} & \textbf{Größen\-ordnungen}\\
\bottomrule
\end{tabular}
\end{center}

\begin{center}
\begin{minipage}{0.88\textwidth}
\kursiv{Die zweite Zeile ist die wichtigste. Die richtige Zahl folgt aus den Anforderungen: Ein Vierpersonenteam braucht zwei Bestandteile, ein Universitätsklinikum vielleicht zwölf. Wer eine Zahl nennt, ohne die Anforderungen zu kennen, rät.}
\end{minipage}
\end{center}
\end{warnung}
\end{merksatz}

\bigskip
\begin{flushright}
\itshape\color{grau}
Ein Programm löst eine Aufgabe.\\
Ein Werk hält ein Versprechen ---\\
über Jahre, über Ausfälle,\\
über alle hinweg, die daran gebaut haben.\\[10pt]
\normalfont\color{grau}\upshape
--- die Bilanz dieses Bandes.\\
\itshape Das Versprechen steht in den Nähten.\\
Nicht im Code.
\end{flushright}
\vspace{6pt}
\noindent Kleppmann\zit{15} ordnet die Ent\-schei\-dun\-gen ver\-teil\-ter Werke; Brewer\zit{4} liefert den Rahmen für die Ab\-wä\-gung; Abbott und Fisher\zit{1} behandeln das Wachs\-tum solcher Werke über Jahre.
